% GENERATED FILE. Edit canonical lessons, then run npm run generate:books.
% canonical-source-hash: fnv1a64:1206c2ea4dd9af78
% canonical-lessons: UR-C123-nadi, UR-C123-mitti, UR-C123-ghas, UR-C123-bij, UR-C123-jar

\chapter{Stream, Earth, Grass, Seed, Root}
\label{ch:ur-stream-earth-grass-seed-root}
\hlchaptermodality{123}

\begin{chapteropening}
\textbf{By the end of this chapter:} \emph{I can say a stream, earth, grass, a seed and a root.}

Complete the last lesson of chapter 123: I can say a stream, earth, grass, a seed and a root.
\end{chapteropening}

\section[\texorpdfstring{\ur{ندی} (nadī)}{ندی (nadī)}]{\ur{ندی} (nadī) — a stream, a small river}

\label{lesson:UR-C123-nadi}



\begin{quote}
Before the new one: say the Urdu for gloves, then the Urdu for a hill.
\end{quote}

\subsection*{You'll want to know: \ur{ندی}}
\textbf{\ur{ندی}} — \emph{nadī} — \textquotedblleft{}a stream, a small river\textquotedblright{}.

\begin{cousinweb}
\textbf{\ur{ندی}} (\emph{nadī}) is Sanskrit \emph{nadī}, \textquotedblleft{}river\textquotedblright{}, taken over whole as a learned loan, a tatsama (CDIAL 6943). CDIAL lists the inherited form as \emph{naī}, \textquotedblleft{}small river\textquotedblright{}.
\end{cousinweb}

\begin{grammarlens}[title={What you've built}]
One more word for this chapter.
\end{grammarlens}

\subsection*{Your turn}
\begin{itemize}
\raggedright
  \item \emph{Say it:} \emph{nadī}
  \item \emph{Say it:} \emph{nadī}, once more
  \item \emph{Say it:} say \emph{nadī}
  \item \emph{Recall:} say the Urdu for gloves, then the Urdu for a hill, then say \emph{nadī} again
\end{itemize}

\begin{tcolorbox}[breakable,colback=teal!4,colframe=teal!35!black,title={Before you move on}]
What does \ur{ندی} mean? (\textquotedblleft{}A stream, a small river\textquotedblright{}.) Say it once more.
\end{tcolorbox}
\section[\texorpdfstring{\ur{مٹی} (miṭṭī)}{مٹی (miṭṭī)}]{\ur{مٹی} (miṭṭī) — earth, soil}

\label{lesson:UR-C123-mitti}



\begin{quote}
Before the new one: say the Urdu for a hill, then the Urdu for a stream.
\end{quote}

\subsection*{You'll want to know: \ur{مٹی}}
\textbf{\ur{مٹی}} — \emph{miṭṭī} — \textquotedblleft{}earth, soil\textquotedblright{}.

\begin{cousinweb}
\textbf{\ur{مٹی}} (\emph{miṭṭī}) is a tadbhava, a word inherited from Sanskrit and reshaped by sound change. It comes from \emph{mṛttikā}, \textquotedblleft{}earth, clay\textquotedblright{} (CDIAL 10286).
\end{cousinweb}

\begin{grammarlens}[title={What you've built}]
One more word for this chapter.
\end{grammarlens}

\subsection*{Your turn}
\begin{itemize}
\raggedright
  \item \emph{Say it:} \emph{miṭṭī}
  \item \emph{Say it:} \emph{miṭṭī}, once more
  \item \emph{Say it:} say \emph{miṭṭī}
  \item \emph{Recall:} say the Urdu for a hill, then the Urdu for a stream, then say \emph{miṭṭī} again
\end{itemize}

\begin{tcolorbox}[breakable,colback=teal!4,colframe=teal!35!black,title={Before you move on}]
What does \ur{مٹی} mean? (\textquotedblleft{}Earth, soil\textquotedblright{}.) Say it once more.
\end{tcolorbox}
\section[\texorpdfstring{\ur{گھاس} (ghās)}{گھاس (ghās)}]{\ur{گھاس} (ghās) — grass}

\label{lesson:UR-C123-ghas}



\begin{quote}
Before the new one: say the Urdu for a stream, then the Urdu for earth.
\end{quote}

\subsection*{You'll want to know: \ur{گھاس}}
\textbf{\ur{گھاس}} — \emph{ghās} — \textquotedblleft{}grass\textquotedblright{}.

\begin{cousinweb}
\textbf{\ur{گھاس}} (\emph{ghās}) is a tadbhava, a word inherited from Sanskrit and reshaped by sound change. It comes from \emph{ghāsa}, \textquotedblleft{}food, pasture grass\textquotedblright{} (CDIAL 4471). It has come down with almost no change.
\end{cousinweb}

\begin{grammarlens}[title={What you've built}]
One more word for this chapter.
\end{grammarlens}

\subsection*{Your turn}
\begin{itemize}
\raggedright
  \item \emph{Say it:} \emph{ghās}
  \item \emph{Say it:} \emph{ghās}, once more
  \item \emph{Say it:} say \emph{ghās}
  \item \emph{Recall:} say the Urdu for a stream, then the Urdu for earth, then say \emph{ghās} again
\end{itemize}

\begin{tcolorbox}[breakable,colback=teal!4,colframe=teal!35!black,title={Before you move on}]
What does \ur{گھاس} mean? (\textquotedblleft{}Grass\textquotedblright{}.) Say it once more.
\end{tcolorbox}
\section[\texorpdfstring{\ur{بیج} (bīj)}{بیج (bīj)}]{\ur{بیج} (bīj) — a seed}

\label{lesson:UR-C123-bij}



\begin{quote}
Before the new one: say the Urdu for earth, then the Urdu for grass.
\end{quote}

\subsection*{You'll want to know: \ur{بیج}}
\textbf{\ur{بیج}} — \emph{bīj} — \textquotedblleft{}a seed\textquotedblright{}.

\begin{cousinweb}
\textbf{\ur{بیج}} (\emph{bīj}) is Sanskrit \emph{bīja}, \textquotedblleft{}seed\textquotedblright{}, taken over whole as a learned loan, a tatsama (CDIAL 9250). CDIAL lists the inherited forms of the same word as \emph{biyā} and \emph{bīā}.
\end{cousinweb}

\begin{grammarlens}[title={What you've built}]
One more word for this chapter.
\end{grammarlens}

\subsection*{Your turn}
\begin{itemize}
\raggedright
  \item \emph{Say it:} \emph{bīj}
  \item \emph{Say it:} \emph{bīj}, once more
  \item \emph{Say it:} say \emph{bīj}
  \item \emph{Recall:} say the Urdu for earth, then the Urdu for grass, then say \emph{bīj} again
\end{itemize}

\begin{tcolorbox}[breakable,colback=teal!4,colframe=teal!35!black,title={Before you move on}]
What does \ur{بیج} mean? (\textquotedblleft{}A seed\textquotedblright{}.) Say it once more.
\end{tcolorbox}
\section[\texorpdfstring{\ur{جڑ} (jaṛ)}{جڑ (jaṛ)}]{\ur{جڑ} (jaṛ) — a root}

\label{lesson:UR-C123-jar}



\begin{quote}
Before the new one: say the Urdu for grass, then the Urdu for a seed.
\end{quote}

\subsection*{You'll want to know: \ur{جڑ}}
\textbf{\ur{جڑ}} — \emph{jaṛ} — \textquotedblleft{}a root\textquotedblright{}.

\begin{cousinweb}
\textbf{\ur{جڑ}} (\emph{jaṛ}) is a tadbhava, a word inherited from Sanskrit and reshaped by sound change. It comes from \emph{jaṭā}, \textquotedblleft{}twisted hair; fibrous root\textquotedblright{} (CDIAL 5086). CDIAL weighs a Dravidian and a Munda source for the word.
\end{cousinweb}

\begin{grammarlens}[title={What you've built}]
One more word for this chapter.
\end{grammarlens}

\subsection*{Your turn}
\begin{itemize}
\raggedright
  \item \emph{Say it:} \emph{jaṛ}
  \item \emph{Say it:} \emph{jaṛ}, once more
  \item \emph{Say it:} say \emph{jaṛ}
  \item \emph{Recall:} say the Urdu for grass, then the Urdu for a seed, then say \emph{jaṛ} again
\end{itemize}

\begin{tcolorbox}[breakable,colback=teal!4,colframe=teal!35!black,title={Before you move on}]
What does \ur{جڑ} mean? (\textquotedblleft{}A root\textquotedblright{}.) Say it once more.
\end{tcolorbox}
